\section{Cell measures and mappings}
\label{sec:m3-cells}
\subsection{Finite cells and weighted support}
\label{def:m3-cell-complex}
\paragraph{Setup.}
Let \(K_d\) be a finite set of oriented \(d\)-cells, with free chain group
\(C_d=\mathbb Z[K_d]\), and let \(B_d:C_d\to C_{d-1}\) be its
integer incidence map. A chain complex requires \(B_{d-1}B_d=0\).
The vertices and edges of \eqref{eq:atomic-h1} specify a graph; additional
cells, their incidences and their admission hypotheses are separate choices.
For occupied support \(S\subseteq K_d\), cell weights \(w_d(c)\ge0\), and
multiplicities \(m(c)\in\mathbb N\), define
\begin{equation}
 \mu_d(S)=\sum_{c\in S}w_d(c),\qquad
 \mu_d(S,m)=\sum_{c\in S}m(c)w_d(c).
 \label{eq:m3-support-measure}
\end{equation}
Cell identity defines support; zero weight does not remove a cell from it.
The signed chain \(z=\sum_c n_c c\), its boundary \(B_dz\), and occupied
support are distinct objects. A signed coefficient cancelling to zero
does not imply that either contributing occurrence was absent.

\paragraph{Derivation.}
For disjoint supports, finite summation in \eqref{eq:m3-support-measure}
is additive. For overlapping supports it gives
\(\mu_d(S\cup T)=\mu_d(S)+\mu_d(T)-\mu_d(S\cap T)\).
An occurrence sum instead counts its declared multiplicities. Neither
identity relates a \(d\)-measure to a \((d-1)\)-measure without additional
weights and geometric hypotheses.

\paragraph{Specialization.}
\label{prop:m3-cell-boundary}
For an ordered product of vertices and oriented intervals, let
\(i_1<\cdots<i_d\) denote the interval coordinates and \(c_{j,\pm}\)
the endpoint faces. Set
\begin{equation}
 \partial c=\sum_{j=1}^{d}(-1)^{j-1}(c_{j,+}-c_{j,-}),
 \qquad \partial^2c=0.
 \label{eq:m3-cubical-boundary}
\end{equation}
Indeed, fixing two distinct interval coordinates in either order produces
the same face. Removing the earlier coordinate changes the later index by
one, so the two incidence products have opposite signs. Each codimension-two
face cancels exactly once. Including all faces gives a cubical chain complex.

\paragraph{Falsification.}
Reversing one face incidence while retaining its endpoint incidences can
make \(\partial^2c\ne0\). Counting \(S\cap T\) twice changes the support
measure whenever its weight is positive. A zero-weight shared face still
participates in \eqref{eq:m3-cubical-boundary}.
\aodprovenance{\eqref{eq:atomic-h1}; \eqref{eq:afc-stokes-used};
\eqref{eq:m3-support-measure}.}

\subsection{Folds, projections and a metric counterexample}
\label{prop:m3-fold-support}
\paragraph{Setup.}
A family \(F_d:C_d(K)\to C_d(L)\) is a chain map precisely when
\begin{equation}
 B_d^L F_d=F_{d-1}B_d^K.
 \label{eq:m3-chain-map}
\end{equation}
Such maps may supplement the retained inner identities and scope maps of
\eqref{eq:fold-outer-closure}. That construction alone does not assert
\eqref{eq:m3-chain-map} or preservation of \eqref{eq:m3-support-measure}.

\paragraph{Derivation.}
Commutation sends cycles to cycles, since \(B^Kz=0\) implies
\(B^LFz=FB^Kz=0\).
Take two adjacent oriented unit squares and reflect the right square onto
the left. The induced signed cell maps commute with boundaries; the two
oriented area contributions cancel. The domain occupies two
squares and the image support occupies one. Retaining two occurrences on
that image would give multiplicity two.

\paragraph{Specialization.}
\label{prop:m3-staircase}
A monotone axial quarter path with coordinate displacements \(r,r\),
\(r>0\), has length \(2r\) when its length is the sum of axial displacements.
Four quarters therefore have length \(8r\) at every subdivision. The Euclidean
circle has circumference \(2\pi r\), so subdivision
with the axial estimator retains the discrepancy \((8-2\pi)r\).
For the fixed unit disk and mesh \(1/n\), define quarter-disk upper and
lower rectangle heights
\(u_i=\lceil\sqrt{n^2-i^2}\rceil\) and
\(\ell_i=\lfloor\sqrt{n^2-(i+1)^2}\rfloor\), \(0\le i<n\).
Monotonicity of the circle graph encloses its area between the two rectangle
sums. Writing each ceiling as its floor plus a number in \([0,1]\), the
floor differences telescope to \(n\). Consequently
\begin{equation}
 A_n^-\le A_{\mathrm{disk}}\le A_n^+,\qquad
 A_n^+-A_n^-=\frac{4}{n^2}\sum_{i=0}^{n-1}(u_i-\ell_i)
 \le \frac{8}{n}.
 \label{eq:m3-disk-enclosure}
\end{equation}
The area enclosure converges; the axial perimeter remains fixed.
Likewise, \(C=2qr\) and \(A=qr^2\) imply \(Cr-2A=0\) for every \(q\);
the residual cannot select a coefficient.

\paragraph{Falsification.}
Identify distinct preimages through their equal projected coordinates:
unique reconstruction fails. Reversing an oriented face without its induced
edge map can violate \eqref{eq:m3-chain-map}. Increasing numerical precision
does not change the axial length \(8r\).
\aodprovenance{\eqref{eq:fold-outer-closure}; \eqref{eq:expose-inner-lineage};
\eqref{eq:m3-support-measure}; \eqref{eq:m3-chain-map}.}

\subsection{Typed mappings and conditional error bounds}
\label{sec:m3-maps}
\label{def:m3-map}
\paragraph{Setup.}
Let \(f:D\subseteq X\to Y\) and \(g:E\subseteq Y\to Z\), with specified
quantity dimensions, units and admissibility conditions. Composition requires
\(f(D)\subseteq E\). A projection may be exact and many-to-one: its fibres
state the lost distinctions. A reconstruction error requires a specified
reconstruction map and a comparison metric; it is not implied by noninjectivity.

\paragraph{Derivation.}
\label{prop:m3-composed-bound}
Suppose \(Y,Z\) are normed spaces, both \(f(D)\) and \(\widetilde f(D)\)
lie in \(E\), and on the relevant domains
\(\|\widetilde f-f\|\le b_f\),
\(\|\widetilde g-g\|\le b_g\), and
\(\|g(y)-g(y')\|\le L_g\|y-y'\|\).
Inserting \(g(\widetilde f(x))\) and applying the triangle inequality gives
\begin{equation}
 \|\widetilde g(\widetilde f(x))-g(f(x))\|
 \le b_g+L_g b_f.
 \label{eq:m3-composed-bound}
\end{equation}
These are deterministic bounds in the specified output units.

\paragraph{Specialization.}
For \(g(y)=1/y\) on a positive interval with lower endpoint \(a>0\),
\(L_g=a^{-2}\) suffices. If admissible intermediate values approach zero,
no such finite uniform bound follows; crossing zero makes the reciprocal
undefined there and its remaining image unbounded.

\paragraph{Falsification.}
An approximation taking an intermediate value outside \(E\) violates the
composition hypotheses even if its error is small. On intervals approaching
zero a fixed reciprocal Lipschitz constant fails.
\aodprovenance{The declared domains and norm inequalities of
\secref{sec:m3-maps}; \eqref{eq:m3-composed-bound}.}

\subsection{Finite joint laws and approximation scope}
\label{prop:m3-finite-law}
\paragraph{Setup.}
Let a joint input \(X\) have finite support \(x_s\), probabilities
\(p_s\ge0\), \(\sum_s p_s=1\), and let every \(x_s\) be in the domain
of \(h\). Shared coordinates refer to the same input variable throughout
a composition. The output moments are
\begin{equation}
 \bar y=\sum_s p_sh(x_s),\qquad
 \Sigma_y=\sum_s p_s(h(x_s)-\bar y)(h(x_s)-\bar y)^T.
 \label{eq:m3-finite-law}
\end{equation}
\paragraph{Derivation.}
Finite expectation gives \eqref{eq:m3-finite-law} directly. For any vector
\(v\), \(v^T\Sigma_yv=\sum_sp_s[v^T(h(x_s)-\bar y)]^2\ge0\), so
singular positive-semidefinite covariance is admissible. If \(h\) is affine,
\(h(x)=Ax+b\), substitution gives \(\Sigma_y=A\Sigma_xA^T\) exactly.
For a differentiable nonlinear map, the corresponding expression at
\(\bar x\) is only its first-order covariance:
\begin{equation}
 \Sigma_{y,\mathrm{lin}}=J(\bar x)\Sigma_xJ(\bar x)^T.
 \label{eq:m3-linear-covariance}
\end{equation}
Adequacy for a specified output requires control of the omitted remainder
or a comparison with the full joint law. Means and covariances alone need
not determine nonlinear output moments.

\paragraph{Specialization.}
Let \(X\) take \(-\epsilon,0,\epsilon\) with probabilities
\(1/4,1/2,1/4\), \(\epsilon>0\). For \(h(X)=X^2\), the linearized
variance at zero vanishes, but the exact mean and variance are
\(\epsilon^2/2\) and \(\epsilon^4/4\). In contrast, identical shared
offsets cancel exactly in \((a+X)-(b+X)=a-b\). A constant output on the
entire supplied support has zero variance even for a nonlinear formula.

\paragraph{Falsification.}
Replacing the shared offset by two independent copies changes the model.
Treating the zero derivative of \(X^2\) as a proof of zero spread contradicts
\eqref{eq:m3-finite-law}.
\aodprovenance{\eqref{eq:m3-finite-law}; \eqref{eq:m3-linear-covariance}.}

\subsection{Observation-map identifiability}
\label{prop:m3-atomic-gain}
\paragraph{Setup.}
Use two recurrent coupled lineages admitted by \eqref{eq:relational-time},
with positive window counts \(N_H,N_C\). Suppose an observation law
has free positive gains \(g_H,g_C\) and assigns equivalent cycles
\(C_H=g_HN_H\), \(C_C=g_CN_C\). Fixing a reference scale fixes \(g_C\)
but supplies no constraint on \(g_H/g_C\).
\paragraph{Derivation.}
For any positive proposed ratio \(R\), choosing
\(g_H=R g_C N_C/N_H\) gives \(C_H/C_C=R\), preserving both native
histories and the reference normalization. The ratio is therefore
underdetermined until an independently specified observation law fixes
the relative gain.
\paragraph{Specialization.}
Under a constant reference frequency \(\nu_C>0\), the elapsed
interval is \(C_C/\nu_C\). This additional observation law does not assign
a universal duration to a native cut. A predicted transition period
\(1/\nu_H\) is a derived transformation, distinct from that elapsed interval.
\paragraph{Falsification.}
Selecting \(g_H\) from an evaluation value of \(R\) makes the resulting
ratio a calibration. Changing evaluation data alone must leave a previously
fixed model output unchanged.
\aodprovenance{\eqref{eq:relational-time}; \eqref{eq:cadence-ratio};
\eqref{eq:complete-wave-state}; \eqref{eq:complete-coupling-state}.}
